\section{Quantifier complexity over arithmetic}
\label{sec:decision-hierarchy-02}

\begin{definition}[Language and standard structure of arithmetic]
\label{def:language-arithmetic}
The \emph{language of arithmetic} is
\[
  \Lang_A=\{0,S,+,\times\},
\]
where \(0\) is a constant, \(S\) is unary, and \(+\) and \(\times\) are
binary function symbols.  The \emph{standard structure of arithmetic} is
\[
  \mathcal N=(\N,0,S,+,\times)
\]
with the usual operations.  These objects supply the syntax and semantics for
every hierarchy classification below.
\end{definition}

For \(n\in\N\), the \emph{numeral} for \(n\) is the closed term
\(\bar n=S^n0\).  For \(\bar a=(a_1,\ldots,a_k)\), write
\(\overline{\bar a}=(\bar a_1,\ldots,\bar a_k)\).  The bar is
metalinguistic: it maps a natural number or finite tuple of natural numbers to
an arithmetic term or tuple of arithmetic terms.  When a bar sits over a
single letter carrying a numeric subscript, as in \(\bar a_1\), it denotes
the numeral of that tuple entry; a bar is never applied to a subscripted
tuple name.

For arithmetic terms \(r,t\), write
\(r<t\) as the metalinguistic abbreviation
\(\exists z\,(r+Sz=t)\), where \(z\) is the least variable absent from
\(r,t\).  If \(x\notin\operatorname{Var}(t)\), the marked bounded quantifiers
abbreviate
\[
\begin{aligned}
  \exists x<t\,\varphi &\defeq \exists x\,(x<t\land\varphi),\\
  \forall x<t\,\varphi &\defeq \forall x\,(x<t\to\varphi).
\end{aligned}
\]
Their satisfaction clauses therefore quantify over exactly the natural
numbers smaller than the value of \(t\).  The hierarchy counts each marked
bounded quantifier as bounded before this eliminable expansion.

Alongside \(<\), write \(r\leq t\) as the metalinguistic abbreviation
\(\exists z\,(z+r=t)\), where \(z\) is the least variable absent from
\(r,t\).  If \(x\notin\operatorname{Var}(t)\), the marked bounded
quantifiers
\[
\begin{aligned}
  \exists x\leq t\,\varphi &\defeq \exists x\,(x\leq t\land\varphi),\\
  \forall x\leq t\,\varphi &\defeq \forall x\,(x\leq t\to\varphi)
\end{aligned}
\]
are counted as bounded exactly as their \(<\)-marked counterparts.  In the
standard structure, \(\mathcal N\entails x\leq t\leftrightarrow x<St\);
provability of this equivalence in \(\mathsf Q\) is not claimed and is not
used below.  The witness \(z\) sits to the \emph{left} of \(+\) in
\(z+r=t\); this choice is what lets the comparison lemmas proved for
\(\leq\) in the incompleteness part rest only on the right-recursion
axioms \(x+0=x\) and \(x+Sy=S(x+y)\) for addition.

\begin{definition}[Syntactic arithmetical hierarchy]
A \emph{bounded arithmetic formula} is generated from atomic
\(\Lang_A\)-formulas by \(\bottom,\neg,\land,\lor,\to\) and marked bounded
quantifiers only.

To classify an arbitrary arithmetic formula, first eliminate implications and
push negations to atomic formulas by the De Morgan and quantifier-duality
rewrites, treating a marked bounded quantifier by the corresponding bounded
duality.  Rename bound variables by the fixed least-fresh-variable order when
needed to avoid capture.  This fixed conversion reaches negation normal form;
it does not move an unbounded quantifier across a bounded quantifier or a
connective.

At level zero, the \(\Sigma^0_0\) formulas and the \(\Pi^0_0\) formulas are
exactly the bounded arithmetic formulas in negation normal form.  Having
defined both classes at level \(n\), define the classes at level \(n+1\)
simultaneously as follows:
\begin{itemize}
  \item the \(\Sigma^0_{n+1}\) formulas in negation normal form form the least
        class containing every \(\Sigma^0_n\) and \(\Pi^0_n\) formula and
        closed under \(\land,\lor\), both marked bounded quantifiers, and
        unbounded existential quantification;
  \item the \(\Pi^0_{n+1}\) formulas in negation normal form form the least
        class containing every \(\Sigma^0_n\) and \(\Pi^0_n\) formula and
        closed under \(\land,\lor\), both marked bounded quantifiers, and
        unbounded universal quantification.
\end{itemize}
An arbitrary formula is a \emph{\(\Sigma^0_n\) formula} or a
\emph{\(\Pi^0_n\) formula} iff its fixed negation-normal-form conversion lies
in the corresponding recursively generated class.  The definition therefore
handles a bounded quantifier outside an unbounded quantifier without an
invalid quantifier exchange.  It depends on syntax and the fixed conversion,
not on whether the sentence happens to be true in \(\mathcal N\).
\end{definition}

\begin{definition}[Definability classes of relations]
Let \(k\in\N\).  A relation \(A\subseteq\N^k\) is
\emph{\(\Sigma^0_n\)-definable} iff there is a \(\Sigma^0_n\) formula
\(\varphi(\bar x)\) such that, for every \(\bar a\in\N^k\),
\[
  \bar a\in A
  \quad\Longleftrightarrow\quad
  \mathcal N\entails\varphi(\overline{\bar a}).
\]
Define \(\Pi^0_n\)-definable analogously, and call \(A\)
\emph{\(\Delta^0_n\)-definable} iff it is both
\(\Sigma^0_n\)-definable and \(\Pi^0_n\)-definable.  Write
\[
\begin{aligned}
 \boldsymbol{\Sigma}^0_n(k)&=\{A\subseteq\N^k:A\text{ is }
       \Sigma^0_n\text{-definable}\},\\
 \boldsymbol{\Pi}^0_n(k)&=\{A\subseteq\N^k:A\text{ is }
       \Pi^0_n\text{-definable}\},\\
 \boldsymbol{\Delta}^0_n(k)&=
       \boldsymbol{\Sigma}^0_n(k)\cap\boldsymbol{\Pi}^0_n(k).
\end{aligned}
\]
For unary relations, the argument \((1)\) may be omitted.  The unbold symbols
classify formulas syntactically; the bold symbols classify relations by
definability in \(\mathcal N\).  This typography records different standing
and prevents equivalence in the standard structure \(\mathcal N\) from
reclassifying every true or
false sentence as bounded.  Both standings are used in the results below.
\end{definition}

The semantic fact that a relation is decidable does not by itself provide an
object-language formula for that relation.  The lemmas and coding
definition below supply the particular formulas required in the hierarchy
arguments.  The coding of computation histories is built in-document from
the numerical pairing map of \cref{ch:strings} together with the G\"odel
beta function; every projection used is defined by a bounded formula,
proved as such here.

